\documentclass[11pt]{article}
\usepackage[margin=1in]{geometry}
\usepackage[T1]{fontenc}
\usepackage[utf8]{inputenc}
\usepackage{lmodern}
\usepackage{graphicx}
\usepackage{booktabs}
\usepackage{array}
\usepackage{amsmath}
\usepackage{hyperref}
\usepackage{parskip}
\usepackage{xcolor}

\hypersetup{
  colorlinks=true,
  linkcolor=blue!50!black,
  urlcolor=blue!50!black
}

\title{Citation Imputation and RDD LLM Tiebreaker Counterfactual}
\author{}
\date{\today}

\begin{document}
\maketitle

\section*{Purpose}
This note studies citation imputation and an RDD-style policy counterfactual
using a ridge model on full-submission SPECTER2 embeddings. The narrow
question is:
\emph{if we keep the number of accepted papers fixed near the human
accept/reject cutoff, does replacing some human accepts with high-scoring LLM
rejects move us toward papers with higher downstream citation potential?}

\section*{1. What Was Done}
The workflow has three layers.

First, observed citations were recovered from OpenAlex through the arXiv match
pipeline. In the broader RDD data build, this recovered citations for
11{,}006 paper-level rows out of 11{,}336 arXiv-linked rows, a 97.1\% match
rate conditional on having an arXiv match.

Second, the imputation stage uses a supervised ridge model on full-paper
SPECTER2 embeddings. The full-text representation uses the
extracted submission body with the \texttt{References} section removed. Each
paper is split into evenly spaced chunks across the full body, each chunk is
embedded with SPECTER2, and the paper embedding is the mean of the chunk
vectors. On the 2018--2020 RDD sample, this covered 2{,}390 of 2{,}391 papers;
the one missing full-text file already had observed citations and no LLM score,
so it does not affect the counterfactual.

Third, we run a volume-matched LLM tiebreaker exercise on the 2018--2020 local
RDD sample. The human rule is the sharp cutoff ``accept if mean reviewer
rating is at or above the year-specific threshold.'' Inside a narrow band
around that threshold, we replace the human rule with an LLM committee-score
rule while forcing the within-band acceptance rate to stay constant. This
isolates \emph{who} gets picked rather than changing \emph{how many} papers are
accepted.

\section*{2. How Good Is Ridge Imputation?}
The 2018--2020 counterfactual sample contains 2{,}391 papers. Of these, 1{,}512
have observed OpenAlex citations, 879 require imputation, 2{,}361 have LLM
committee scores, and 2{,}390 have full-text embeddings.

The ridge model is validated on the observed-paper subset: 1{,}511 papers with
observed OpenAlex citations and full-text embeddings. On that set, ridge
achieves Pearson correlation 0.321 and Spearman correlation
0.299 in $\log(1+\text{citations})$, with mean absolute error 1.081 in log
space and 67.1 citations in levels. Ridge predictions are used only where
citations are missing; observed OpenAlex counts remain untouched.

For the 879 unmatched papers that require imputation, the ridge predictions
have mean 30.1 citations, median 27.1, and 90th percentile 50.8. These should
still be read as a rough ranking surrogate rather than literal paper-level
citation forecasts.

\begin{figure}[p]
  \centering
  \includegraphics[width=0.78\textwidth]{figures/ridge_pred_vs_actual.png}
  \caption{Observed-paper validation for the full-text ridge imputer. The model
  preserves useful ordering signal in log citations, but still compresses the
  realized outcome variance, so the imputed values are best treated as a
  quality-ranking surrogate rather than literal paper-level forecasts.}
\end{figure}

\section*{3. Counterfactual with Ridge-Filled Citations}
The main counterfactual now uses
\[
\text{cites\_filled} \sim \text{accepted} + \text{mean\_rating} \mid \text{year\_topic},
\]
where \texttt{year\_topic} is a coarse topic bucket from full-text embeddings
crossed with year. In this model, acceptance carries a coefficient of 0.622,
which corresponds to an 86\% increase in expected citations conditional on
paper-quality controls.

We then reran the volume-matched LLM tiebreaker inside three bands around the
human cutoff: $\delta \in \{0.25, 0.50, 0.75\}$. For each band we keep the
within-band acceptance volume fixed and compare the human rule to the LLM rule.

\begin{center}
\small
\begin{tabular}{>{\raggedright}p{1.25cm}rrrr}
\toprule
Band $\delta$ & Flips each way & Swap gain / flip & Accept gain / flip & Net policy gain \\
\midrule
0.25 & 38  & +150.5 & +55.6 & +1160.4 \\
0.50 & 223 & +30.0  & +36.3 & +467.7 \\
0.75 & 278 & +27.2  & +36.7 & -22.4 \\
\bottomrule
\end{tabular}
\end{center}

The interpretation is stable. At narrow bands, the LLM tiebreaker appears to
pick substantially stronger papers than the human sharp rule leaves out. At the
main operating point, $\delta = 0.50$, 223 rejects are flipped in and 223
accepts are flipped out. The average citation potential of the flipped-in set
is 59.9, compared with 29.9 for the flipped-out set, a gain of 30.0 per swap.
The Poisson model implies that the flipped-in papers would gain another 36.3
citations each if they had been accepted. The overall net policy gain is
+467.7 predicted citations.

The design lesson is clear: the LLM is most useful as a \emph{local
tiebreaker}, not as a broad override rule.

Composition remains an important caveat. At the main $\delta = 0.50$ band,
31.4\% of the flipped-in papers rely on imputed rather than observed citations.
So the sign of the result is informative, but part of the magnitude still
depends on a surrogate outcome model.

\begin{figure}[p]
  \centering
  \includegraphics[width=0.86\textwidth]{figures/fig_swap_and_gain_by_band.png}
  \caption{Main policy results by band width under the full-text ridge
  imputer. Narrow-band gains remain clearly positive; the broad-band policy is
  close to break-even rather than strongly negative.}
\end{figure}

\begin{figure}[p]
  \centering
  \includegraphics[width=0.98\textwidth]{figures/fig_band_scatter.png}
  \caption{Borderline papers at the main band width $\delta = 0.50$. The
  dashed vertical line is the human cutoff; the dashed horizontal line is the
  volume-matched LLM cutoff. Shape marks whether the citation outcome is
  observed or ridge-imputed.}
\end{figure}

\section*{4. Takeaways}
The citation imputer in this analysis is a supervised ridge model on
full-submission embeddings. The imputed values should still be read as a rough
quality-ranking surrogate rather than literal paper-level forecasts, but they
are strong enough to support the local policy exercise.

The policy conclusion is limited but positive. The evidence supports a narrow
claim: \emph{an LLM committee may be useful as a volume-matched tiebreaker
among borderline papers}. The results do not support using the LLM as a broad
override rule away from the acceptance margin.

\end{document}
